\chapter{The Implementation Treasury: Theory \& Code}
\label{app:treasury}

This appendix provides a unified reference for the mathematical formulations and 
PyTorch implementations discussed throughout the book. Each theoretical concept 
is paired with a self-contained code snippet for immediate practical application.

% -----------------------------------------------------------------------------
\section{Shared Foundation}
% -----------------------------------------------------------------------------

All code snippets in this treasury assume the following standard imports and 
PyTorch environment:

\begin{lstlisting}[style=pythonstyle, caption={Shared imports}, label=lst:imports]
import math
import torch
import torch.nn as nn
import torch.nn.functional as F
from typing import Optional, List
\end{lstlisting}

% -----------------------------------------------------------------------------
\section{Optimization Objectives}
\label{app:objectives}
\label{app:pipeline}
% -----------------------------------------------------------------------------
% -----------------------------------------------------------------------------

A production LLM is produced by five sequential optimisation stages.

\subsection{Stage 1 --- Pre-training}
\label{form:ntp}

\textbf{Objective} (causal language modelling):
\[
  \theta_1^* = \arg\min_\theta\; \mathcal{L}_\mathrm{PT}(\theta) = -\frac{1}{|\mathcal{D}_\mathrm{PT}|}\sum_{(x_1,\ldots,x_T)\in\mathcal{D}_\mathrm{PT}} \sum_{t=1}^{T} \log p_\theta(x_t \mid x_{<t})
\]

\begin{lstlisting}[style=pythonstyle, caption={Causal LM loss with token masking}, label=lst:ntp_loss]
def causal_lm_loss(logits: torch.Tensor, labels: torch.Tensor) -> torch.Tensor:
    """Computes cross-entropy loss by shifting logits and labels for NTP.
    Labels with -100 are ignored (standard for SFT masking).
    """
    # Shift: Predict token at t from tokens 0 to t-1
    shift_logits = logits[:, :-1, :].contiguous().view(-1, logits.size(-1))
    shift_labels = labels[:, 1:].contiguous().view(-1)
    return F.cross_entropy(shift_logits, shift_labels, ignore_index=-100)
\end{lstlisting}

\subsection{Perplexity (PPL)}
\label{form:ppl}

\[
  \mathrm{PPL} = \exp(\mathcal{L}_\mathrm{NTP})
\]
Perplexity measures the model's confidence in its predictions; it represents the average branching factor of the distribution.

\subsection{Stage 3 --- Preference Optimisation}

\noindent\textbf{DPO} (Direct Preference Optimization):
\[
  \mathcal{L}_\mathrm{DPO}(\theta) = -\mathbb{E}_{(x,y_w,y_l)}\!\left[ \log\sigma\!\left( \beta\log\frac{\pi_\theta(y_w|x)}{\pi_\mathrm{ref}(y_w|x)} - \beta\log\frac{\pi_\theta(y_l|x)}{\pi_\mathrm{ref}(y_l|x)} \right) \right]
\]

\noindent\textbf{SimPO} (Simple Preference Optimization):
\[
  \mathcal{L}_\mathrm{SimPO} = -\mathbb{E}\!\left[\log\sigma\!\left( \frac{\beta}{|y_w|}{\log\pi_\theta(y_w|x)} - \frac{\beta}{|y_l|}{\log\pi_\theta(y_l|x)} - \gamma \right)\right]
\]

\noindent\textbf{GRPO} (Group Relative Policy Optimisation):
\[
  \hat{A}_i = \frac{r_i - \mu_r}{\sigma_r}
\]

% -----------------------------------------------------------------------------
\section{Attention Mechanisms}
\label{app:attention}
% -----------------------------------------------------------------------------

\subsection{Scaled Dot-Product Attention}
\label{form:sdpa}

\[
  \mathrm{Attention}(Q, K, V) = \mathrm{softmax}\!\left(\frac{Q K^\top}{\sqrt{d_k}} + M\right) V
\]

\begin{lstlisting}[style=pythonstyle, caption={Scaled dot-product attention}, label=lst:sdpa]
def scaled_dot_product_attention(
    Q: torch.Tensor,   # (Batch, Heads, QueryLen, Dim)
    K: torch.Tensor,   # (Batch, Heads, KeyLen, Dim)
    V: torch.Tensor,   # (Batch, Heads, KeyLen, ValDim)
    causal: bool = False,
) -> torch.Tensor:
    """Computes scaled dot-product attention with optional causal masking."""
    d_k = Q.size(-1)
    scores = Q @ K.transpose(-2, -1) / math.sqrt(d_k)
    
    if causal:
        T, S = scores.shape[-2], scores.shape[-1]
        mask = torch.triu(torch.ones(T, S, device=Q.device), diagonal=1).bool()
        scores = scores.masked_fill(mask, float('-inf'))
        
    attn_weights = F.softmax(scores, dim=-1)
    return attn_weights @ V
\end{lstlisting}

\subsection{Multi-Head Attention (MHA)}
\label{form:mha}

\[
  \mathrm{head}_h = \mathrm{Attention}(Q W_h^Q, K W_h^K, V W_h^V), \quad
  \mathrm{MHA}(Q, K, V) = \mathrm{Concat}(\mathrm{head}_1, \dots, \mathrm{head}_H) W^O
\]

\begin{lstlisting}[style=pythonstyle, caption={Multi-head attention}, label=lst:mha]
class MultiHeadAttention(nn.Module):
    def __init__(self, d_model: int, n_heads: int):
        super().__init__()
        self.h, self.d_k = n_heads, d_model // n_heads
        self.Wq = nn.Linear(d_model, d_model, bias=False)
        self.Wk = nn.Linear(d_model, d_model, bias=False)
        self.Wv = nn.Linear(d_model, d_model, bias=False)
        self.Wo = nn.Linear(d_model, d_model, bias=False)

    def _split_into_heads(self, x: torch.Tensor) -> torch.Tensor:
        B, T, _ = x.shape
        return x.view(B, T, self.h, self.d_k).transpose(1, 2)

    def forward(self, x: torch.Tensor, causal: bool = True) -> torch.Tensor:
        Q = self._split_into_heads(self.Wq(x))
        K = self._split_into_heads(self.Wk(x))
        V = self._split_into_heads(self.Wv(x))
        out = scaled_dot_product_attention(Q, K, V, causal=causal)
        B, H, T, dk = out.shape
        out = out.transpose(1, 2).reshape(B, T, H * dk)
        return self.Wo(out)
\end{lstlisting}

\subsection{Grouped Query Attention (GQA)}
\label{form:gqa}

KV groups shared by $H/G$ query heads. $G=1$ is MQA; $G=H$ is MHA.

\begin{lstlisting}[style=pythonstyle, caption={Grouped query attention}, label=lst:gqa]
class GroupedQueryAttention(nn.Module):
    def __init__(self, d_model: int, n_heads: int, n_kv_groups: int):
        super().__init__()
        self.h, self.g = n_heads, n_kv_groups
        self.d_k = d_model // n_heads
        self.Wq = nn.Linear(d_model, n_heads * self.d_k, bias=False)
        self.Wk = nn.Linear(d_model, n_kv_groups * self.d_k, bias=False)
        self.Wv = nn.Linear(d_model, n_kv_groups * self.d_k, bias=False)
        self.Wo = nn.Linear(d_model, d_model, bias=False)

    def forward(self, x: torch.Tensor, causal: bool = True) -> torch.Tensor:
        B, T, _ = x.shape
        q = self.Wq(x).view(B, T, self.h, self.d_k).transpose(1, 2)
        k = self.Wk(x).view(B, T, self.g, self.d_k).transpose(1, 2)
        v = self.Wv(x).view(B, T, self.g, self.d_k).transpose(1, 2)
        reps = self.h // self.g
        k = k.repeat_interleave(reps, dim=1)
        v = v.repeat_interleave(reps, dim=1)
        out = scaled_dot_product_attention(q, k, v, causal=causal)
        out = out.transpose(1, 2).reshape(B, T, self.h * self.d_k)
        return self.Wo(out)
\end{lstlisting}

\subsection{Multi-Query Attention (MQA)}
\label{form:mqa}

Shares a single KV pair across all query heads. Special case of GQA where $G=1$.

\begin{lstlisting}[style=pythonstyle, caption={Multi-query attention}, label=lst:mqa]
class MultiQueryAttention(nn.Module):
    def __init__(self, d_model: int, n_heads: int):
        super().__init__()
        self.h, self.d_k = n_heads, d_model // n_heads
        self.Wq = nn.Linear(d_model, d_model, bias=False)
        self.Wk = nn.Linear(d_model, self.d_k, bias=False)
        self.Wv = nn.Linear(d_model, self.d_k, bias=False)
        self.Wo = nn.Linear(d_model, d_model, bias=False)

    def forward(self, x: torch.Tensor, causal: bool = True) -> torch.Tensor:
        B, T, _ = x.shape
        q = self.Wq(x).view(B, T, self.h, self.d_k).transpose(1, 2)
        k = self.Wk(x).view(B, T, 1, self.d_k).transpose(1, 2)
        v = self.Wv(x).view(B, T, 1, self.d_k).transpose(1, 2)
        out = scaled_dot_product_attention(q, k, v, causal=causal)
        out = out.transpose(1, 2).reshape(B, T, self.h * self.d_k)
        return self.Wo(out)
\end{lstlisting}

\subsection{Multi-head Latent Attention (MLA)}
\label{form:mla}

\[
  c_t^{KV} = W^{DKV} h_t, \quad [k_t^C; v_t^C] = W^{UKV} c_t^{KV}
\]

\begin{lstlisting}[style=pythonstyle, caption={Multi-head latent attention (MLA)}, label=lst:mla]
class MLA(nn.Module):
    def __init__(self, d_model: int, n_heads: int, d_latent: int, d_v_latent: int):
        super().__init__()
        self.h, self.d_k = n_heads, d_model // n_heads
        self.W_dk_c = nn.Linear(d_model, d_latent, bias=False)
        self.W_uk = nn.Linear(d_latent, n_heads * self.d_k, bias=False)
        self.W_uv = nn.Linear(d_latent, n_heads * self.d_k, bias=False)
        self.Wo = nn.Linear(d_model, d_model, bias=False)

    def forward(self, x: torch.Tensor, causal: bool = True) -> torch.Tensor:
        B, T, _ = x.shape
        c_kv = self.W_dk_c(x)
        k = self.W_uk(c_kv).view(B, T, self.h, self.d_k).transpose(1, 2)
        v = self.W_uv(c_kv).view(B, T, self.h, self.d_k).transpose(1, 2)
        q = nn.Linear(x.shape[-1], self.h * self.d_k, device=x.device)(x).view(B, T, self.h, self.d_k).transpose(1, 2)
        out = scaled_dot_product_attention(q, k, v, causal=causal)
        out = out.transpose(1, 2).reshape(B, T, self.h * self.d_k)
        return self.Wo(out)
\end{lstlisting}

% -----------------------------------------------------------------------------
\section{Positional Encodings}
\label{app:pos}
% -----------------------------------------------------------------------------

\subsection{Rotary Position Embedding (RoPE)}
\label{form:rope}

\[
  R_{\theta_i, t} = \begin{pmatrix} \cos(t\theta_i) & -\sin(t\theta_i) \\ \sin(t\theta_i) & \cos(t\theta_i) \end{pmatrix}, \quad \theta_i = 10000^{-2(i-1)/d}
\]

\begin{lstlisting}[style=pythonstyle, caption={Rotary positional embeddings (RoPE)}, label=lst:rope]
def apply_rotary_emb(x: torch.Tensor, freqs: torch.Tensor) -> torch.Tensor:
    """Applies complex rotary embeddings to the last dimension of x."""
    x_complex = torch.view_as_complex(x.float().reshape(*x.shape[:-1], -1, 2))
    freqs = freqs.view(1, 1, x.shape[-2], -1)
    x_rotated = x_complex * torch.exp(1j * freqs)
    return torch.view_as_real(x_rotated).flatten(-2).type_as(x)

def precompute_freqs(dim: int, seq_len: int, theta: float = 10000.0) -> torch.Tensor:
    """Precomputes the rotary frequency tensor for a given context length."""
    freqs = 1.0 / (theta ** (torch.arange(0, dim, 2)[: (dim // 2)].float() / dim))
    t = torch.arange(seq_len)
    return torch.outer(t, freqs).float()
\end{lstlisting}

\subsection{ALiBi (Attention with Linear Biases)}
\label{form:alibi}

\[
  \mathrm{score}_{ij} = \frac{q_i \cdot k_j}{\sqrt{d_k}} - m_h \cdot |i - j|
\]

\begin{lstlisting}[style=pythonstyle, caption={ALiBi bias matrix}, label=lst:alibi]
def alibi_bias(n_heads: int, seq_len: int) -> torch.Tensor:
    """Generates the linear penalty matrix for ALiBi attention."""
    slopes = 2 ** (-8 * torch.arange(1, n_heads + 1) / n_heads)
    queries = torch.arange(seq_len).unsqueeze(0)
    keys    = torch.arange(seq_len).unsqueeze(1)
    pos     = queries - keys
    bias = -slopes.view(-1, 1, 1) * pos.abs().unsqueeze(0).float()
    return bias.tril()
\end{lstlisting}

% -----------------------------------------------------------------------------
\section{Normalization and Activation}
\label{app:norm}
% -----------------------------------------------------------------------------

\subsection{RMSNorm}
\label{form:rmsnorm}

\[
  \mathrm{RMSNorm}(x) = \gamma \odot \frac{x}{\|x\|_\mathrm{RMS} + \epsilon}, \quad \|x\|_\mathrm{RMS} = \sqrt{\frac{1}{d}\sum_ix_i^2}
\]

\begin{lstlisting}[style=pythonstyle, caption={RMSNorm}, label=lst:rmsnorm]
class RMSNorm(nn.Module):
    """Root Mean Square Layer Normalization (standard in Llama 2/3)."""
    def __init__(self, d_model: int, eps: float = 1e-8):
        super().__init__()
        self.eps, self.gamma = eps, nn.Parameter(torch.ones(d_model))

    def forward(self, x: torch.Tensor) -> torch.Tensor:
        rms = x.pow(2).mean(-1, keepdim=True).add(self.eps).sqrt()
        return self.gamma * x / rms
\end{lstlisting}

\subsection{SwiGLU Feed-Forward Network}
\label{form:swiglu}

\[
  \mathrm{SwiGLU}(x, W, V, W_2) = (x W \odot \mathrm{Swish}(x V)) W_2
\]

\begin{lstlisting}[style=pythonstyle, caption={SwiGLU FFN}, label=lst:swiglu]
class SwiGLU(nn.Module):
    """Gated Linear Unit with Swish activation (standard in LLM FFNs)."""
    def __init__(self, d_model: int, d_ff: int):
        super().__init__()
        self.W1 = nn.Linear(d_model, d_ff, bias=False)
        self.W2 = nn.Linear(d_model, d_ff, bias=False)
        self.W3 = nn.Linear(d_ff, d_model, bias=False)

    def forward(self, x: torch.Tensor) -> torch.Tensor:
        return self.W3(F.silu(self.W1(x)) * self.W2(x))
\end{lstlisting}

% -----------------------------------------------------------------------------
\section{Optimization}
\label{app:optim}
% -----------------------------------------------------------------------------

\subsection{AdamW}
\label{form:adamw}

\[
  \theta_t = \theta_{t-1} - \eta \left(\frac{\hat{m}_t}{\sqrt{\hat{v}_t}+\epsilon} + \lambda \theta_{t-1}\right)
\]

\begin{lstlisting}[style=pythonstyle, caption={AdamW update step (manual)}, label=lst:adamw]
def adamw_step(
    params: list[torch.Tensor], grads: list[torch.Tensor], 
    m: list[torch.Tensor], v: list[torch.Tensor], t: int,
    lr: float = 3e-4, b1: float = 0.9, b2: float = 0.95, 
    eps: float = 1e-8, wd: float = 0.1
) -> int:
    """Manual implementation of the AdamW update rule with decoupled weight decay."""
    t += 1
    for p, g, m_i, v_i in zip(params, grads, m, v):
        m_i.mul_(b1).add_(g, alpha=1 - b1)
        v_i.mul_(b2).addcmul_(g, g, value=1 - b2)
        m_hat = m_i / (1 - b1 ** t)
        v_hat = v_i / (1 - b2 ** t)
        p.mul_(1 - lr * wd)
        p.addcdiv_(m_hat, v_hat.sqrt().add_(eps), value=-lr)
    return t
\end{lstlisting}

\subsection{Muon Optimizer}
\label{form:muon}

\[
  X_{new} = 0.5 X (3I - XX^\top)
\]

\begin{lstlisting}[style=pythonstyle, caption={Muon update step}, label=lst:muon]
@torch.no_grad()
def muon_step(p, g, m, lr=0.02, momentum=0.95, n_iters=5):
    """Orthogonalizing optimizer for the 2026 frontier."""
    m.mul_(momentum).add_(g)
    g = g.add(m, alpha=momentum)
    X = g.view(g.size(0), -1)
    X /= X.norm().clamp(min=1e-7)
    for _ in range(n_iters):
        X = 0.5 * X @ (3.0 * torch.eye(X.size(0), device=X.device) - X @ X.T)
    p.add_(X.view_as(p), alpha=-lr)
\end{lstlisting}

\subsection{Lion Optimizer}
\label{form:lion}

\[
  c_t = \text{sign}(\beta_1 m_{t-1} + (1-\beta_1) g_t), \quad m_t = \beta_2 m_{t-1} + (1-\beta_2) g_t
\]
Lion uses the `sign` operation to determine update direction, providing better generalization in specific LLM regimes.

\subsection{Cosine Schedule with Warmup}
\label{form:cosine}

\begin{lstlisting}[style=pythonstyle, caption={Cosine LR schedule with warmup}, label=lst:cosine]
def cosine_lr(step: int, total_steps: int, lr_max: float, lr_min: float, warmup_steps: int) -> float:
    if step < warmup_steps:
        return lr_max * step / warmup_steps
    progress = (step - warmup_steps) / max(1, total_steps - warmup_steps)
    return lr_min + 0.5 * (lr_max - lr_min) * (1 + math.cos(math.pi * progress))
\end{lstlisting}

% -----------------------------------------------------------------------------
\section{Parameter-Efficient Fine-Tuning}
\label{app:peft}
% -----------------------------------------------------------------------------

\subsection{LoRA (Low-Rank Adaptation)}
\label{form:lora}

\[
  h = W_0 x + \frac{\alpha}{r} B A x
\]

\begin{lstlisting}[style=pythonstyle, caption={LoRA linear layer}, label=lst:lora]
class LoRALinear(nn.Module):
    def __init__(self, base: nn.Linear, rank: int = 16, alpha: float = 32.0):
        super().__init__()
        d_out, d_in = base.weight.shape
        self.base, self.scale = base, alpha / rank
        self.base.weight.requires_grad_(False)
        self.A = nn.Parameter(torch.randn(rank, d_in) * 0.01)
        self.B = nn.Parameter(torch.zeros(d_out, rank))

    def forward(self, x):
        return self.base(x) + (x @ self.A.T @ self.B.T) * self.scale
\end{lstlisting}

\subsection{DoRA (Weight-Decomposed Low-Rank Adaptation)}
\label{form:dora}

\[
  W = m \frac{W_0 + BA}{\|W_0 + BA\|_F}
\]

\subsection{PPO (Proximal Policy Optimization)}
\label{form:ppo}

\[
  \mathcal{L}_\mathrm{CLIP} = \hat{\mathbb{E}}_t [ \min(r_t(\theta) \hat{A}_t, \text{clip}(r_t(\theta), 1-\epsilon, 1+\epsilon) \hat{A}_t) ]
\]

\begin{lstlisting}[style=pythonstyle, caption={PPO clipped objective}, label=lst:ppo]
def ppo_loss(policy_logits, old_logits, actions, advantages, eps=0.2):
    ratio = (policy_logits - old_logits).exp()
    surr1 = ratio * advantages
    surr2 = torch.clamp(ratio, 1-eps, 1+eps) * advantages
    return -torch.min(surr1, surr2).mean()
\end{lstlisting}

% -----------------------------------------------------------------------------
\section{Alignment Objectives}
\label{app:alignment}
% -----------------------------------------------------------------------------

\subsection{DPO (Direct Preference Optimization)}
\label{form:dpo}

\[
  \mathcal{L}_\mathrm{DPO}(\theta) = -\mathbb{E}_{(x,y_w,y_l)}\!\left[ \log\sigma\!\left( \beta\log\frac{\pi_\theta(y_w|x)}{\pi_\mathrm{ref}(y_w|x)} - \beta\log\frac{\pi_\theta(y_l|x)}{\pi_\mathrm{ref}(y_l|x)} \right) \right]
\]

\begin{lstlisting}[style=pythonstyle, caption={Direct Preference Optimization loss}, label=lst:dpo]
def dpo_loss(policy, ref, x_w, x_l, beta=0.1):
    """Computes the DPO loss for a pair of responses."""
    pi_w, pi_l = log_prob_seq(policy, x_w, x_w), log_prob_seq(policy, x_l, x_l)
    ref_w, ref_l = log_prob_seq(ref, x_w, x_w), log_prob_seq(ref, x_l, x_l)
    reward_margin = beta * ((pi_w - ref_w) - (pi_l - ref_l))
    return -F.logsigmoid(reward_margin).mean()
\end{lstlisting}

\subsection{SimPO (Simple Preference Optimization)}
\label{form:simpo}

\[
  \mathcal{L}_\mathrm{SimPO} = -\mathbb{E}\!\left[\log\sigma\!\left( \frac{\beta}{|y_w|}{\log\pi_\theta(y_w|x)} - \frac{\beta}{|y_l|}{\log\pi_\theta(y_l|x)} - \gamma \right)\right]
\]

\begin{lstlisting}[style=pythonstyle, caption={SimPO loss}, label=lst:simpo]
def simpo_loss(model, x_w, x_l, beta=2.0, gamma=0.5):
    """Length-normalized preference optimization without a reference model."""
    r_w, r_l = avg_log_prob(model, x_w, x_w), avg_log_prob(model, x_l, x_l)
    return -F.logsigmoid(beta * (r_w - r_l) - gamma).mean()
\end{lstlisting}

\subsection{GRPO (Group Relative Policy Optimisation)}
\label{form:grpo}

\[
  \hat{A}_i = \frac{r_i - \mu_r}{\sigma_r}, \quad \mathcal{L}_\mathrm{GRPO}(\theta) = -\mathbb{E}\!\left[\sum_{i=1}^{G}\sum_t \dots \right]
\]

\begin{lstlisting}[style=pythonstyle, caption={GRPO update step}, label=lst:grpo]
def grpo_step(
    policy: nn.Module, 
    ref: nn.Module, 
    prompt_ids: torch.Tensor, 
    reward_fn: callable,
    G: int = 8, 
    beta: float = 0.001, 
    eps: float = 0.2
) -> torch.Tensor:
    """Group Relative Policy Optimization update step."""
    with torch.no_grad():
        completions = [policy.generate(prompt_ids) for _ in range(G)]
        rewards     = reward_fn(completions)
        old_logp    = torch.stack([log_prob_seq(policy, c, c) for c in completions])

    advantage = (rewards - rewards.mean()) / (rewards.std() + 1e-8)

    total_batch_loss = 0.0
    for i, comp in enumerate(completions):
        curr_logp = log_prob_seq(policy, comp, comp)
        ratio     = (curr_logp - old_logp[i]).exp()
        
        surr1     = ratio * advantage[i]
        surr2     = torch.clamp(ratio, 1 - eps, 1 + eps) * advantage[i]
        pg_loss   = -torch.min(surr1, surr2)

        kl_div    = curr_logp - log_prob_seq(ref, comp, comp)
        total_batch_loss += pg_loss + beta * kl_div

    return total_batch_loss / G
\end{lstlisting}

\subsection{KTO (Kahneman-Tversky Optimization)}
\label{form:kto}

\[
  \mathcal{L}_\mathrm{KTO} = -\mathbb{E}[w \cdot \log \sigma(\beta(r(x,y) - \text{ref}))]
\]

\begin{lstlisting}[style=pythonstyle, caption={KTO loss}, label=lst:kto]
def kto_loss(model, ref, input_ids, labels, beta=0.1):
    """Alignment from binary feedback using Kahneman-Tversky utility."""
    logp = log_prob_seq(model, input_ids, labels)
    logp_ref = log_prob_seq(ref, input_ids, labels)
    return -F.logsigmoid(beta * (logp - logp_ref)).mean()
\end{lstlisting}

% -----------------------------------------------------------------------------
\section{Inference and Distillation}
\label{app:inference_distill}
\label{app:distill}
% -----------------------------------------------------------------------------

\subsection{Knowledge Distillation (Forward KL)}
\label{form:fwd_kd}

\[
  \mathcal{L}_\mathrm{fwd} = D_\mathrm{KL}(p_\mathrm{teacher} \| p_\mathrm{student})
\]

\begin{lstlisting}[style=pythonstyle, caption={Forward KL distillation}, label=lst:fwd_kd]
def fwd_kl_distill(s_logits, t_logits, temp=2.0):
    p = F.softmax(t_logits / temp, dim=-1)
    log_q = F.log_softmax(s_logits / temp, dim=-1)
    return F.kl_div(log_q, p, reduction="batchmean") * (temp**2)
\end{lstlisting}

\subsection{MiniLLM (Reverse KL)}
\label{form:rev_kd}

\[
  \mathcal{L}_\mathrm{rev} = D_\mathrm{KL}(p_\mathrm{student} \| p_\mathrm{teacher})
\]

\subsection{KV Cache Management}
\label{form:kvcache}

The memory footprint (in bytes) of the KV cache is:
\[
  M_{KV} = 2 \times L \times T \times H \times d_k \times \text{bytes\_per\_param}
\]
where $L$ is layers, $T$ is sequence length, and $H$ is heads.

\subsection{Quantization Scaling}
\label{form:quant}

Linear quantization to $n$ bits (e.g., INT8):
\[
  x_q = \text{round}\!\left( \frac{x}{\Delta} \right), \quad \Delta = \frac{\max(|x|)}{2^{n-1}-1}
\]

\subsection{Speculative Decoding}
\label{form:spec_dec}

Accept token $x_t$ with probability $\alpha(x_t) = \min(1, p(x_t)/q(x_t))$.

\begin{lstlisting}[style=pythonstyle, caption={Speculative decoding loop}, label=lst:spec_dec]
@torch.no_grad()
def speculative_decode(
    target: nn.Module, 
    draft: nn.Module, 
    input_ids: torch.Tensor, 
    max_new: int = 128, 
    k: int = 4
) -> torch.Tensor:
    """Accelerated inference via draft-target asymmetric sampling."""
    ids = input_ids.clone()
    while ids.shape[1] < input_ids.shape[1] + max_new:
        draft_ids, draft_probs = ids.clone(), []
        for _ in range(k):
            d_logits = draft(draft_ids).logits[:, -1]
            p_draft  = F.softmax(d_logits, dim=-1)
            tok      = torch.multinomial(p_draft, 1)
            draft_probs.append(p_draft[0, tok[0, 0]].item())
            draft_ids = torch.cat([draft_ids, tok], dim=1)

        t_logits = target(draft_ids).logits
        accepted = 0
        for j in range(k):
            p_tgt   = F.softmax(t_logits[0, ids.shape[1]+j-1], dim=-1)
            tok_j   = draft_ids[0, ids.shape[1] + j].item()
            alpha   = min(1.0, p_tgt[tok_j].item() / (draft_probs[j] + 1e-9))
            if torch.rand(1).item() < alpha:
                accepted += 1
            else:
                break
        
        ids = torch.cat([ids, draft_ids[:, ids.shape[1]:ids.shape[1]+accepted+1]], dim=1)
    return ids
\end{lstlisting}

% -----------------------------------------------------------------------------
\section{Model Merging}
\label{app:merging}
% -----------------------------------------------------------------------------

\subsection{Task Vectors}
\label{form:task_vec}

\[
  \tau = \theta_{fine-tuned} - \theta_{base}
\]

\subsection{SLERP (Spherical Linear Interpolation)}
\label{form:slerp}

\[
  \mathrm{SLERP}(\theta_A, \theta_B, t) = \frac{\sin((1-t)\Omega)}{\sin\Omega}\theta_A + \frac{\sin(t\Omega)}{\sin\Omega}\theta_B
\]

\begin{lstlisting}[style=pythonstyle, caption={SLERP weight interpolation}, label=lst:slerp]
def slerp(w_a: torch.Tensor, w_b: torch.Tensor, t: float) -> torch.Tensor:
    """Interpolates between weights along an arc for better distribution preservation."""
    a, b = w_a.flat().float(), w_b.flat().float()
    cos_omega = (a / a.norm()).dot(b / b.norm()).clamp(-1, 1)
    omega = cos_omega.acos()
    if omega.abs() < 1e-6: return (1-t)*w_a + t*w_b
    return (math.sin((1-t)*omega)*a + math.sin(t*omega)*b) / math.sin(omega)
\end{lstlisting}

\subsection{DARE (Drop And REscale)}
\label{form:dare}

\[
  \theta_{merge} = \theta_{base} + \frac{1}{1-p} \sum_{i=1}^n \text{mask}_i \odot \tau_i
\]
where $p$ is the dropout probability.

\subsection{TIES-Merging}
\label{form:ties}

Three-step conflict resolution: Trim, Elect sign, and Disjoint merge.

\begin{lstlisting}[style=pythonstyle, caption={Task Vector}, label=lst:task_vec]
def task_vector(base_sd: dict, ft_sd: dict) -> dict:
    """Computes the delta vector: theta_diff = theta_ft - theta_base."""
    return {k: ft_sd[k].float() - base_sd[k].float() for k in base_sd}
\end{lstlisting}

\begin{lstlisting}[style=pythonstyle, caption={TIES-Merging}, label=lst:ties]
def ties_merge(base_sd: dict, ft_sds: list[dict], lambdas: list[float], density: float = 0.7) -> dict:
    """Trim, Elect, and Merge conflict resolution."""
    taus = [task_vector(base_sd, ft) for ft in ft_sds]
    trimmed = []
    # ... (Trimming and logic)
    return merged

\end{lstlisting}

% -----------------------------------------------------------------------------
\section{Continual Learning --- EWC}
\label{app:continual}
% -----------------------------------------------------------------------------

\subsection{Elastic Weight Consolidation (EWC)}
\label{form:ewc}

\begin{lstlisting}[style=pythonstyle, caption={Elastic Weight Consolidation}, label=lst:ewc]
class EWC:
    def __init__(self, model, dataset, lam=1000.0):
        self.lam, self.params = lam, {n: p.clone().detach() for n, p in model.named_parameters()}
        self.fisher = self._compute_fisher(model, dataset)
    # ... (Penalty and Fisher logic)
\end{lstlisting}

% -----------------------------------------------------------------------------
\section{Final Assembly --- TinyGPT}
\label{app:tinygpt}
% -----------------------------------------------------------------------------

A minimal, production-ready decoder architecture incorporating the treasury's components.

\begin{lstlisting}[style=pythonstyle, caption={Complete TinyGPT architecture}, label=lst:tinygpt]
class DecoderLayer(nn.Module):
    def __init__(self, d_model: int = 512, n_heads: int = 8, n_kv: int = 2, d_ff: int = 2048):
        super().__init__()
        self.norm1, self.norm2 = RMSNorm(d_model), RMSNorm(d_model)
        self.attn, self.ffn = GroupedQueryAttention(d_model, n_heads, n_kv), SwiGLU(d_model, d_ff)

    def forward(self, x: torch.Tensor, freqs: torch.Tensor) -> torch.Tensor:
        x = x + self.attn(self.norm1(x))
        return x + self.ffn(self.norm2(x))

class TinyGPT(nn.Module):
    def __init__(self, vocab=32000, d=512, heads=8, kv=2, ff=2048, layers=6, max_seq=2048):
        super().__init__()
        self.embed   = nn.Embedding(vocab, d)
        self.layers  = nn.ModuleList([DecoderLayer(d, heads, kv, ff) for _ in range(layers)])
        self.norm, self.head = RMSNorm(d), nn.Linear(d, vocab, bias=False)
        self.head.weight = self.embed.weight 
        self.register_buffer("freqs", precompute_freqs(d // heads, max_seq))

    def forward(self, idx: torch.Tensor) -> torch.Tensor:
        x = self.embed(idx)
        for layer in self.layers: x = layer(x, self.freqs)
        return self.head(self.norm(x))
\end{lstlisting}
